\subsection{The free group of rank two}\label{sec:free}

The last and most important example in the gallery. Every earlier group arrived with its elements
already in hand: rotations, permutations, remainders, pairs of integers. Here the \emph{set} is
defined by a condition and the \emph{rule} is a two-step recipe, so even closure has to be argued,
and associativity becomes a genuine theorem.

\emph{The alphabet.} Four symbols: $a$, $b$, $a\inv$, $b\inv$. At this stage $a\inv$ is nothing but
a symbol; the notation is chosen because it will turn out to be an inverse once the group exists.

\emph{Words.} A \defn{word} is any finite string of these four symbols, including the
\defn{empty word}, which has no letters at all. A word is \defn{reduced} if it never has a symbol
standing next to its partner, that is, no $aa\inv$, $a\inv a$, $bb\inv$ or $b\inv b$ anywhere in it.

\emph{The set.} $F_2=\{\text{reduced words in }a,b\}$.

\emph{The rule.} Write one word after the other, then delete forbidden pairs until none is left.

In the figures each letter is a box, red for the $a$ family and blue for the $b$ family, and a
shaded pair is the one about to be deleted.

\sfigs{0.86}{A word, a word that is not reduced, and the reduced word left after deleting.}{\figXam}

\subsubsection*{Closure}

Two things need checking here, and both are easy to overlook. First, the deleting \emph{stops}:
each deletion removes two boxes, so a string of length $k$ can be cut at most $k/2$ times. Second,
what it stops at is a \emph{reduced} word, because that is precisely the condition for no deletion
being available. So the answer lands back in $F_2$.

There is also a useful picture of where the cancelling happens. Both factors were already reduced,
so the only forbidden pair in the concatenation is at the \emph{seam}, and the deleting unzips
outwards from there.

\sfigs{0.86}{Multiplying $aba\inv$ by $ab$: only the seam can create a forbidden pair.}{\figXan}

\subsubsection*{Identity}

The empty word. Writing nothing in front of or behind a word $w$ changes no letter and creates no
forbidden pair, so $1w=w1=w$.

\sfigs{0.86}{The identity of $F_2$ is the word with no letters.}{\figXao}

\subsubsection*{Inverses}

Reverse the word and flip every letter, turning $a$ into $a\inv$, $a\inv$ into $a$, and likewise for
$b$. This is the socks-and-shoes rule again: to undo ``socks, then shoes'' you take the shoes off
first. The two strings then annihilate from the middle outwards.

\sfigs{0.86}{$w=abb$ against $w\inv=b\inv b\inv a\inv$: the middle pair goes first, then the next, then the
last, leaving the empty word.}{\figXap}

\subsubsection*{Associativity}

This is the axiom that costs something. $(uv)w$ and $u(vw)$ are two different \emph{orders of
deleting} inside the one long string $uvw$, and there is no reason on the face of it why different
orders should reach the same reduced word. Take $u=ab$, $v=b\inv a\inv$, $w=ab$.

\sfigs{0.86}{$(uv)w$: the product $uv$ collapses completely, and the answer is $ab$.}{\figXaq}

\sfigs{0.86}{$u(vw)$: now $vw$ collapses completely, and the answer is again $ab$. The two routes deleted
different pairs in a different order.}{\figXar}

\begin{thm}\label{thm:unique}
Every word in $a,b,a\inv,b\inv$ reduces to exactly one reduced word, whatever order the deletions
are performed in.
\end{thm}

\begin{proof}[Proof sketch]
Induct on the length of the word. If two deletions are available at once, they either involve four
distinct positions, in which case performing one leaves the other still available and the two orders
lead to the same shorter word; or they overlap in a position, in which case the overlap forces a
pattern such as $x\,x\inv\,x$, and deleting either pair leaves the same single letter $x$. Either
way any two first steps can be followed by further deletions reaching a common word, and the
induction hypothesis applies to it.
\end{proof}

\sfig{The word $a\,a\inv\,a\,b\,b\inv\,a$ has three places where a deletion is available. Whichever
you choose, and whatever you do next, the process ends at the same reduced word $aa$.}{\figreduce}

\begin{why}
\emph{Why Theorem~\ref{thm:unique} gives associativity.} Reduce the single string $uvw$ in any order
at all. Both $(uv)w$ and $u(vw)$ are such reductions, since each is obtained by deleting pairs
inside $uvw$, merely in a prescribed order. By the theorem they end at the same reduced word, so
they are equal, which is (G3). \emph{This is the price of defining a group by strings instead of by
functions:} for the cube, for $S_n$ and for $D_n$ associativity came free from composition of
functions.
\end{why}

\sfigs{0.86}{Multiplying in $F_2$: concatenate, then reduce.}{\figXas}

\sfigs{0.86}{Checking an inverse: $aba\inv b$ against $b\inv ab\inv a\inv$ collapses from the middle
outwards to the empty word.}{\figXat}

\subsubsection*{The extra properties}

$\{a,b\}$ generates $F_2$ by construction, since every element \emph{is} a word in $a$ and $b$ and
their inverses. The group is infinite; more precisely there are $4\cdot 3^{\,k-1}$ reduced words of
length $k$, because after the first letter each new letter may be anything except the partner of the
one before, so the number of elements grows exponentially with length.

It is not abelian: $ab$ and $ba$ are different reduced words, hence different elements. In fact no
element except the identity has finite order, since a nonempty reduced word never collapses when
repeated.

The word \defn{free} means precisely that the only relations holding in $F_2$ are the forced ones
$xx\inv=1$. Nothing else is ever true there. Two consequences: when you build a homomorphism out of
$F_2$ you may send $a$ and $b$ absolutely anywhere (Section~\ref{sec:homs}), and every group
whatever can be obtained from a free group by imposing extra relations
(Section~\ref{sec:presentations}).

\begin{pitfall}
Compare $\Z^2$ of Section~\ref{sec:Z2}: the same two generators, one extra rule $ab=ba$, and a
completely different group. In $\Z^2$ there are about $k^2$ elements of length $k$; in $F_2$ there
are about $3^k$. Adding one relation changed polynomial growth into exponential growth.
\end{pitfall}

\sfig{The words $ab$ and $ba$ as paths, and why they end at different places.}{\figabba}

\sfigs{0.80}{\label{fig:tree}The object whose symmetry group is $F_2$, drawn the same way as every
other picture in this chapter: one dot per group element, and one coloured arrow per generator.
Start at the dot marked $1$ and read a word left to right, taking a red edge rightwards for $a$ and
backwards along a red edge for $a^{-1}$, a blue edge upwards for $b$ and backwards for $b^{-1}$. The
shaded path is the word $ab$: right, then up. Edges are drawn shorter at each step outwards, so that
the whole tree fits; every dot really has four edges at it. Because a reduced word never doubles
back, distinct reduced words end at distinct dots, which is why $F_2$ is infinite and why the
picture is a tree rather than a grid.}{\figTree}

\begin{tryit}
Reduce $b\,a\,a\inv\,b\inv\,b\,a$ in two different orders and confirm that both give $ba$. Write
down the inverse of $a\inv bab\inv$ and check the collapse. How many reduced words have length $3$?
Finally, decide which axiom would fail if the set were taken to be \emph{all} words, not just the
reduced ones, with plain concatenation as the rule.
\end{tryit}
